% !TEX root = ../main.tex

% 学士学位论文大摘要
% 相较于前置摘要，大摘要需要更完整地呈现研究背景、动机、方法、实验与结论。
% 本文件在前置中英文摘要之外，补充了对三章核心工作和系统验证结果的展开介绍。

% 2025 版写作规范要求末尾保留英文大摘要，这里保留中文内容但不参与排版。
\iffalse
\begin{digest}[zh]

\subsection*{研究背景与动机}

视觉目标导航是具身智能中连接感知与决策的核心任务。给定机器人当前视觉观察与目标图像，
策略需要在无高精度地图、无全局定位先验的前提下输出可执行的局部运动，
从而完成从起点到目标视角的自主移动。近年来，以 NoMaD 为代表的通用导航模型
将扩散策略引入视觉导航，把未来局部轨迹建模为条件分布，通过去噪采样生成多模态候选，
相比回归单步速度的端到端策略具有显著更强的表达能力。但现有工作的闭环验证主要集中于
轮式平台（如 LoCoBot、TurtleBot、Jackal），面向四足机器人的部署链路仍缺乏系统研究。
四足平台在步态抖动、控制频率、身体姿态耦合等方面与轮式平台差异明显，
这为高层扩散策略与中低层控制接口的协同提出了新的问题。

\subsection*{研究目标与技术路线}

本文以 NoMaD 为高层视觉导航核心，以云深处 Lite3 四足机器人为目标部署平台，
回答三个面向部署的问题：第一，NoMaD 的视觉编码器、距离预测头与扩散策略头在代码
实现上如何耦合，推理层有哪些优化空间；第二，是否可以在不重新训练模型的前提下，
通过采样加速、条件引导、测试时搜索与视觉编码器替换等手段显著提升策略可用性；
第三，这些优化能否整合进面向 Lite3 的分层闭环系统，并完成 MuJoCo 仿真与
Jetson Orin 真机部署的完整验证。技术路线沿"算法层优化—系统层集成—部署层验证"
三级递进展开，对应论文第三、四、五章。

\subsection*{主要工作与创新点}

\textbf{算法层}：建立了统一的离线统计实验链路，对 DDIM 推理加速、Classifier-Free
Guidance (CFG) 引导增强、Test-Time Scaling (TTS) 推理时搜索与视觉编码器升级进行了
系统对比。DDIM 实验中 \texttt{ddim\_2} 配置相对 DDPM-10 基线实现约 4.84 倍加速且
轨迹质量无显著退化；CFG 实验通过 score function 修改分析给出引导尺度选择准则，
指出其有效范围有限、宜作为可调参数；TTS 实验提出将总收益分解为"评分筛选"
（近乎零时延代价）与"候选池扩容"（时延线性增长但边际递减）两部分，TTS-8 在三因素
联合排序中以最小代价取得综合最优，为推理预算分配提供了新的诠释；视觉编码器升级
采用两阶段框架比较 EfficientNet-B0、ResNet-50、DINOv2-Small、ConvNeXt-Tiny 的
离线指标，为后续部署筛选出候选组合。

\textbf{系统层}：将统一推理模块、状态机、导航主机、中层控制映射与平台接口整合为
分层闭环系统，支持 MuJoCo 仿真与真机两种后端。状态机由 idle / stand / navigate /
explore / recovery / estop / completed / failed 组成，通过统一接口与平台解耦；
导航主机支持批量计划（JSON）和交互式任务两种模式，允许在同一会话内完成拍照、建图、
导航、探索等组合任务。MuJoCo 闭环实验中，Lite3 在 easy/medium/hard 三张地图上均
成功完成定点导航、多目标导航、自由探索任务，所有目标均被到达。

\textbf{部署层}：针对 Jetson Orin 与 Lite3 上位机链路，形成了分阶段联调流程、
多级安全机制与图像记录规范。Orin 阶段一（独立调试）验证了相机采集、模型加载、
推理基准与端到端流水线；在桥接层对速度命令加入 NaN/Inf 守卫、幅值限幅与
连接健康监测，任意推理异常自动走安全急停；状态机层加入全局摔倒保护，使
idle/stand 等状态也不再遗漏摔倒判定。

\subsection*{实验与结论}

离线实验表明，DDIM-2 + 可选 CFG + 可配置 TTS 是当前最值得进入闭环验证的组合；
MuJoCo 闭环表明分层系统能够稳定完成四足导航与探索任务，多目标任务中的任务切换
不会破坏视觉上下文；Orin 阶段一的独立测试验证了真机流水线的基础可用性。
本文将 NoMaD 从轮式平台的实验框架扩展为面向四足平台的可分析、可仿真、可部署的
研究链路，并沉淀出一套可复用的推理优化实验协议、闭环系统组织方式以及真机部署
安全护栏清单，为后续基于扩散策略的足式机器人导航研究提供了参考基础。

\end{digest}
\fi

\begin{digest}[en]

\subsection*{Background and Motivation}

Visual goal navigation is a core task in embodied intelligence that connects perception
with decision-making. Given the current visual observation and a goal image, the policy
must output executable local motion without relying on high-precision maps or global
localization, thereby achieving autonomous movement from the current viewpoint to the
goal viewpoint. Recently, general navigation models such as NoMaD have introduced
diffusion policies into visual navigation, modelling the future local trajectory as a
conditional distribution and generating multi-modal candidates through denoising
sampling, with significantly stronger expressive power than regression-based end-to-end
policies. However, closed-loop validation of these models has so far concentrated on
wheeled platforms (LoCoBot, TurtleBot, Jackal), while deployment pipelines targeting
legged robots remain underexplored. Legged platforms differ substantially from wheeled
ones in gait-induced jitter, control-loop frequency, and body-posture coupling, which
poses new challenges for the collaboration between high-level diffusion policies and
mid-/low-level control interfaces.

\subsection*{Objectives and Approach}

This thesis takes NoMaD as the high-level navigation core and the DeepRobotics Lite3
quadruped as the target deployment platform, and addresses three deployment-oriented
questions: (i) how the visual encoder, distance-prediction head, and diffusion policy
head are coupled in the NoMaD implementation and where inference-level optimizations are
possible; (ii) whether, without retraining, sampling acceleration, conditional guidance,
test-time search, and encoder replacement can meaningfully improve policy usability; and
(iii) whether these optimizations can be integrated into a hierarchical closed-loop
system for Lite3 and validated through MuJoCo simulation and Jetson Orin real-robot
deployment. The technical roadmap proceeds through algorithm-level optimization,
system-level integration, and deployment-level validation, corresponding respectively
to Chapters 3, 4, and 5.

\subsection*{Main Contributions}

\textbf{Algorithm layer.} A unified offline statistical pipeline is built to compare
DDIM inference acceleration, Classifier-Free Guidance (CFG) enhancement, Test-Time
Scaling (TTS) search, and visual encoder upgrades. DDIM-2 attains about $4.84\times$
speedup over the DDPM-10 baseline without notable trajectory-quality degradation. The
CFG analysis, derived from score-function modification, establishes a guidance-scale
selection rule and shows that CFG is best used as a tunable knob rather than a default.
The TTS analysis decomposes total gain into a nearly zero-latency \emph{scoring-based
selection} component and a linearly priced but rapidly diminishing \emph{candidate-pool
expansion} component; TTS-8 wins the three-factor joint ranking at minimal cost,
providing a novel interpretation of inference-budget allocation. Encoder upgrades use
a two-stage framework to compare EfficientNet-B0, ResNet-50, DINOv2-Small, and
ConvNeXt-Tiny on offline metrics.

\textbf{System layer.} A hierarchical closed-loop system---comprising a unified
inference module, state machine, navigation host, mid-level control mapping, and
platform interface---supports both MuJoCo and real-robot backends. The state machine
(idle / stand / navigate / explore / recovery / estop / completed / failed) is decoupled
from platforms through a uniform interface. The navigation host supports both batched
JSON plans and interactive commands, allowing capture, topomap building, navigation,
and exploration to be composed within one session. In MuJoCo closed-loop experiments,
Lite3 successfully completes point navigation, multi-goal navigation, and free
exploration on easy / medium / hard maps, reaching every goal.

\textbf{Deployment layer.} For the Jetson Orin host and the Lite3 upper-computer link,
the thesis develops a staged commissioning procedure, multi-level safety mechanisms,
and an image-recording convention. Orin Stage 1 (standalone commissioning) validates
camera capture, model loading, inference benchmarking, and the end-to-end pipeline. The
bridging layer adds NaN/Inf guards, amplitude clamping, and connection-health checks to
the velocity commands, routing any inference anomaly to a safe stop. A global fall check
is added at the state-machine level so that idle / stand states no longer miss
fallen-robot situations.

\subsection*{Experiments and Conclusions}

Offline results indicate that DDIM-2 plus optional CFG and configurable TTS is the most
promising configuration to bring into closed-loop validation. MuJoCo experiments confirm
that the hierarchical system can robustly complete quadruped navigation and exploration
tasks, and that multi-goal task switching preserves visual context. Orin Stage-1
validation demonstrates the baseline feasibility of the real-robot pipeline. The thesis
extends NoMaD from a wheeled-platform framework into an analyzable, simulatable, and
deployable research pipeline for quadruped platforms, and distills a reusable
inference-optimization protocol, a closed-loop system organization, and a real-robot
safety-guard checklist---providing a practical basis for subsequent research on
diffusion-policy-based visual navigation for legged robots.

\end{digest}
